\documentclass[11pt]{article}
\usepackage[utf8]{inputenc}
\usepackage[margin=1in]{geometry}
\usepackage{amsmath}
\usepackage{graphicx}
\usepackage{booktabs}
\usepackage{hyperref}
\usepackage{natbib}
\usepackage{setspace}
\onehalfspacing

\title{Sex-Specific Resilience of Neocortex to Food Restriction}
\author{Author A$^{1}$, Author B$^{1,2}$, Author C$^{2}$ \\
\small $^{1}$Department of Neuroscience, University of Research \\
\small $^{2}$Department of Physiology, University of Research}
\date{}

\begin{document}
\maketitle

\begin{abstract}
The brain consumes approximately 20\% of the body's caloric intake, yet whether the neocortex contributes to energy-saving strategies during food scarcity is poorly understood, particularly with respect to sex differences. Previous work demonstrated that food restriction reduces energy usage and orientation selectivity in primary visual cortex (V1) of male mice, mediated by decreased serum leptin. Here we tested whether females show similar cortical energy-saving adaptations. We found that females required 40\% greater food restriction (39\% vs 27\% of ad libitum intake) to achieve equivalent 15\% weight loss. Serum leptin decreased approximately 3-fold in food-restricted males ($p < 0.0001$) but showed no significant change in females ($p = 0.32$). RNA-seq analysis of V1 tissue revealed that cellular energy pathways (AMPK, PPAR$\alpha$, mTOR, oxidative phosphorylation) were significantly modulated by food restriction in males but not females. AMPK Thr172 phosphorylation increased in male V1 ($p = 0.022$) but not female V1 ($p = 0.11$). PPARalpha DNA binding activity increased in males ($p = 0.013$) but not females ($p = 0.99$). Two-photon ATP imaging demonstrated that visual stimulation-evoked ATP usage was reduced by food restriction in males but not females. Orientation selectivity in V1 was decreased by food restriction in males but not significantly in females. Bayes factor analysis confirmed more robust effects in males across all parameters. These results demonstrate that female neocortex is resilient to food restriction, maintaining normal energy usage and sensory coding despite caloric deficit, likely through preservation of serum leptin levels.
\end{abstract}

\section{Introduction}

When faced with food scarcity, mammals activate energy-conserving mechanisms across multiple organ systems to prolong survival \citep{wang2006,rosenbaum2010}. These responses include reduced metabolic rate, decreased thermogenesis, and altered hormonal signaling. The adipokine leptin serves as a key signal linking peripheral energy status to central nervous system function: leptin levels decrease proportionally with fat mass during caloric restriction, triggering adaptive responses including reduced energy expenditure and increased appetite \citep{friedman2019,ahima1996}.

The brain is the most energy-expensive organ in mammals, consuming approximately 20\% of total caloric intake in humans despite representing only 2\% of body mass \citep{harris2012}. Within the brain, the neocortex accounts for a substantial fraction of energy consumption, with much of this energy devoted to synaptic transmission and the maintenance of membrane potentials during sensory processing \citep{attwell2001}. Whether the neocortex participates in energy-saving strategies during food restriction has only recently been investigated.

Previous work demonstrated that chronic food restriction reduces energy usage in primary visual cortex (V1) of male mice, as measured by two-photon ATP imaging, and simultaneously decreases orientation selectivity, a key metric of visual cortical function \citep{padamsey2022}. These effects were mediated by decreased serum leptin, as exogenous leptin administration reversed both the energy reduction and the orientation selectivity decline. However, these experiments were conducted exclusively in male mice, leaving open the question of whether females show similar cortical adaptations.

Sex differences in metabolic responses to food restriction are well established in peripheral tissues. Females tend to maintain fat reserves more tenaciously than males, requiring greater caloric restriction to achieve equivalent weight loss \citep{valle2005}. Sex differences in leptin regulation are also well documented, with females generally having higher leptin levels at equivalent body mass index and potentially different leptin sensitivity \citep{rosenbaum2005}. Whether these peripheral sex differences translate to differential cortical responses to food restriction has not been examined.

Understanding sex-specific cortical responses to food restriction has practical implications for neuroscience research. Food restriction is widely used as a motivational tool in behavioral experiments, and its effects on cortical function could confound experimental results in a sex-dependent manner. More broadly, understanding whether the neocortex contributes to energy conservation in both sexes informs our understanding of how the brain prioritizes resources during metabolic stress.

\section{Methods}

\subsection{Animals and Food Restriction Protocol}

Adult C57BL/6 mice (7--9 weeks, both sexes) were assigned to control (CTR, ad libitum feeding) or food-restricted (FR) groups. FR mice received reduced daily food portions calibrated to achieve 85\% of free-feeding weight over 2--3 weeks. All physiological measurements were performed after re-feeding to satiety, isolating long-term caloric restriction effects from acute hunger states.

\begin{table}[h]
\centering
\caption{Animal group characteristics.}
\label{tab:animals}
\begin{tabular}{llll}
\toprule
Group & Weight [95\% CI] & Diet & Restriction Level \\
\midrule
CTR Male ($n = 8$--17) & 27.82 g [26.84--28.79] & Ad libitum & --- \\
FR Male ($n = 8$--19) & Reduced to 85\% & Restricted & 27\% \\
CTR Female ($n = 7$--23) & 23.21 g [22.35--24.07] & Ad libitum & --- \\
FR Female ($n = 6$--11) & Reduced to 85\% & Restricted & 39\% \\
\bottomrule
\end{tabular}
\end{table}

\subsection{Serum Leptin ELISA}

Blood was collected via cardiac puncture at sacrifice. Serum leptin was measured using a mouse leptin ELISA kit according to the manufacturer's protocol.

\subsection{Western Blot for AMPK Phosphorylation}

V1 tissue was dissected, lysed, and subjected to SDS-PAGE and Western blotting for phospho-AMPK (Thr172) and total AMPK. Phosphorylation levels were normalized to total AMPK protein.

\subsection{PPAR$\alpha$ DNA Binding Assay}

Nuclear extracts from V1 tissue were assayed for PPAR$\alpha$ DNA binding activity using a transcription factor ELISA kit, normalized to total nuclear protein.

\subsection{RNA Sequencing}

V1 tissue from CTR and FR mice of both sexes was processed for RNA-seq using standard library preparation and Illumina sequencing. Differentially expressed genes (DEGs) were identified using DESeq2 with an adjusted $p$-value threshold of 0.1. Pathway analysis was performed using Gene Set Enrichment Analysis (GSEA).

\subsection{Two-Photon ATP Imaging}

ATeam1.03YEMK transgenic mice expressing the ATP FRET sensor were used for two-photon imaging of ATP dynamics in V1 layer 2/3 during natural scene presentation. ATP synthesis inhibitors (oligomycin + 2-deoxyglucose) were applied to isolate ATP usage from production. Changes in FRET ratio during visual stimulation reflect ATP consumption.

\subsection{Orientation Selectivity}

Two-photon calcium imaging of V1 L2/3 neurons was performed during drifting grating stimulation (12 directions, 30$^{\circ}$ apart). The orientation selectivity index (OSI) was calculated for each responsive neuron.

\subsection{Statistical Analysis}

Group comparisons used two-sample $t$-tests. Sex $\times$ diet interactions were assessed with two-way ANOVA. Bayes factors (BF$_{10}$) were computed using default priors to supplement frequentist statistics, providing evidence for or against the alternative hypothesis \citep{rouder2009}.

\section{Results}

\subsection{Females Require Greater Food Restriction for Equivalent Weight Loss}

Both sexes achieved the target 15\% weight loss, but females required substantially more restriction (Table~\ref{tab:food}). FR males consumed 27\% less food than CTR males, while FR females consumed 39\% less---a 40\% greater restriction level. Both CTR male vs FR male ($t_{25} = 4.81$, $p < 0.0001$) and CTR female vs FR female ($t_{25} = 6.59$, $p < 0.0001$) showed significant differences in food intake. CTR males and CTR females consumed similar amounts ($t_{25} = 0.013$, $p = 0.99$), while the difference between FR males and FR females trended toward significance ($t_{25} = 1.93$, $p = 0.07$).

\begin{table}[h]
\centering
\caption{Food intake and weight comparisons.}
\label{tab:food}
\begin{tabular}{lccc}
\toprule
Comparison & $t$-statistic & df & $p$-value \\
\midrule
CTR male vs FR male (intake) & 4.81 & 25 & $< 0.0001$ \\
CTR female vs FR female (intake) & 6.59 & 25 & $< 0.0001$ \\
CTR male vs CTR female (intake) & 0.013 & 25 & 0.99 \\
FR male vs FR female (intake) & 1.93 & 25 & 0.07 \\
CTR male vs FR male (weight) & 11.36 & 25 & $< 0.0001$ \\
CTR female vs FR female (weight) & 8.36 & 25 & $< 0.0001$ \\
\bottomrule
\end{tabular}
\end{table}

\subsection{Serum Leptin Maintained in Females Despite Food Restriction}

Food restriction produced a dramatic approximately 3-fold decrease in serum leptin in males ($t_{66} = 4.58$, $p < 0.0001$; $n = 17$ CTR, 19 FR). In striking contrast, female leptin levels showed no significant change ($t_{66} = 1.00$, $p = 0.32$; $n = 23$ CTR, 11 FR). This sex difference in leptin regulation is consistent with the greater metabolic resilience of females and provides a potential mechanistic explanation for the downstream cortical effects.

\subsection{Molecular Energy Pathways Modulated in Males Only}

AMPK Thr172 phosphorylation, a master sensor of cellular energy status, was significantly increased in V1 of FR males ($t_{39} = 2.28$, $p = 0.022$; $n = 11$ CTR, 15 FR) but not FR females ($t_{39} = 0.64$, $p = 0.11$; $n = 11$ CTR, 6 FR). PPAR$\alpha$ DNA binding activity, a downstream effector of energy-sensing pathways, was strongly increased in FR males ($t_{30} = 4.81$, $p = 0.013$; $n = 5$ CTR, 11 FR) but showed no change in FR females ($t_{30} = 0.0016$, $p = 0.99$; $n = 9$ CTR, 9 FR).

\subsection{RNA-Seq Reveals Sex-Specific Transcriptomic Responses}

RNA-seq analysis of V1 tissue revealed a striking sex difference in the transcriptomic response to food restriction (Table~\ref{tab:rnaseq}). In males, food restriction produced a significant number of DEGs (adjusted $p < 0.1$), with GSEA identifying significant modulation of AMPK, mTOR, PPAR$\alpha$, and oxidative phosphorylation pathways. In females, far fewer DEGs were identified, and none of the energy-regulating pathways reached significance. The overlap between male and female DEGs was minimal, as visualized in Venn diagram analysis.

\begin{table}[h]
\centering
\caption{RNA-seq pathway modulation by food restriction in V1.}
\label{tab:rnaseq}
\begin{tabular}{lcc}
\toprule
Pathway & Males: FR Effect & Females: FR Effect \\
\midrule
AMPK signaling & Significantly modulated & Not significant \\
mTOR signaling & Significantly modulated & Not significant \\
PPAR$\alpha$ signaling & Significantly modulated & Not significant \\
Oxidative phosphorylation & Significantly modulated & Not significant \\
\bottomrule
\end{tabular}
\end{table}

\subsection{ATP Usage Reduced in Males but Not Females}

Two-photon imaging of ATP dynamics during natural scene presentation revealed that food restriction significantly decreased visual stimulation-evoked ATP usage in male V1 L2/3 but had no significant effect in females. In the absence of visual stimulation (dark conditions), ATP usage was unaffected by food restriction in either sex, confirming that the effect is specific to activity-dependent energy consumption during sensory processing.

\subsection{Orientation Selectivity Decreased in Males but Not Females}

Consistent with the energy reduction, orientation selectivity in V1 L2/3 neurons was significantly decreased by food restriction in males but showed no significant change in females. Population coding precision followed the same pattern, with FR males showing reduced precision relative to CTR males and no significant difference in females.

\subsection{Bayes Factor Analysis Confirms Sex-Specific Effects}

Bayes factor analysis provided converging evidence across all measures (Table~\ref{tab:bayes}). For each measure, males showed moderate to strong evidence for an effect of food restriction (BF$_{10} > 3$), while females showed weak evidence (BF$_{10} < 3$), consistent with the absence of significant frequentist effects.

\begin{table}[h]
\centering
\caption{Bayes factor analysis summary.}
\label{tab:bayes}
\begin{tabular}{lcccc}
\toprule
Parameter & Males BF$_{10}$ & Females BF$_{10}$ & Males Effect & Females Effect \\
\midrule
Serum leptin & $> 10$ & $< 3$ & Large & Small \\
AMPK phosphorylation & Moderate--strong & Weak & Moderate & Small \\
PPAR$\alpha$ activity & Strong & Weak/absent & Large & Negligible \\
ATP usage & Moderate--strong & Weak & Moderate & Small \\
Orientation selectivity & Moderate & Weak & Moderate & Small \\
\bottomrule
\end{tabular}
\end{table}

\subsection{Pupil Diameter Control}

Pupil diameter during visual stimulation was not significantly affected by sex or food restriction ($t_{18} = 0.82$, $p = 0.42$), ruling out differences in arousal or attention as confounds for the cortical measurements.

\section{Discussion}

Our results demonstrate a striking sex difference in the neocortical response to food restriction. While male mouse V1 undergoes a coordinated energy-saving response---with activation of AMPK and PPAR$\alpha$ pathways, reduced ATP usage during visual processing, and decreased orientation selectivity---female V1 appears largely resilient to the same metabolic challenge. This resilience is associated with the maintenance of serum leptin levels in food-restricted females, despite a greater degree of caloric restriction.

The preservation of leptin signaling in females likely reflects the well-documented sex difference in adipose tissue biology. Females maintain higher body fat percentages and preferentially preserve fat stores during caloric restriction \citep{valle2005}. Since leptin is secreted proportionally to fat mass, preserved fat stores in females would maintain leptin at levels sufficient to prevent the activation of cortical energy-saving mechanisms. This interpretation is consistent with previous work showing that exogenous leptin can reverse food restriction effects in male V1 \citep{padamsey2022}.

The practical implications for neuroscience research are important. Food restriction is widely used to motivate behavioral performance in mice, and our results indicate that this manipulation differentially affects cortical function depending on sex. Male mice undergoing food restriction may have altered sensory processing that could confound behavioral results, while females may be less affected. These sex-specific effects should be considered when designing experiments that combine food restriction with sensory or cognitive testing.

The convergence of molecular (AMPK, PPAR$\alpha$, RNA-seq), metabolic (ATP imaging), and functional (orientation selectivity) measures provides a comprehensive picture of the male-specific cortical response. The cascade appears to proceed from decreased leptin to activated AMPK signaling to reduced ATP consumption to decreased neural coding precision. Each step in this cascade is specifically absent in females, consistent with the upstream maintenance of leptin as the key protective factor.

Several limitations should be noted. First, the two-photon imaging experiments used different mouse lines (ATeam transgenic for ATP, GCaMP for calcium), and line-specific effects cannot be entirely excluded. Second, the FR protocol achieved equivalent weight loss but with different restriction levels, and some differences might relate to the degree rather than the sex specificity of restriction. Third, we did not test whether exogenous leptin reduction in females would render their cortex sensitive to food restriction, which would provide stronger causal evidence for the leptin-mediated mechanism.

\section{Conclusion}

Female mouse neocortex demonstrates resilience to food restriction, maintaining normal energy usage and sensory coding despite significant caloric deficit. This resilience is associated with preserved serum leptin levels and the absence of energy-sensing pathway activation in visual cortex. These findings reveal a sex-specific cortical adaptation to metabolic stress and have important implications for the design and interpretation of neuroscience experiments using food-restricted animals.

\bibliographystyle{plainnat}
\bibliography{references}

\end{document}
